% CV — Amine Mokhtar. Build with:  pdflatex cv_amine_mokhtar.tex
\documentclass[10pt,a4paper]{article}
\usepackage[T1]{fontenc}
\usepackage[utf8]{inputenc}
% Optional, if your TeX installation has French support:
% \usepackage[french]{babel}
\usepackage{lmodern}
\usepackage[a4paper,top=1.25cm,bottom=1.2cm,left=1.4cm,right=1.4cm]{geometry}
\usepackage{enumitem}
\usepackage{titlesec}
\usepackage{microtype}
\usepackage[hidelinks]{hyperref}

\hypersetup{
  pdftitle={CV - Amine Mokhtar},
  pdfauthor={Amine Mokhtar},
  pdfsubject={Élève ingénieur Cloud \& DevOps}
}

\pagestyle{empty}
\setlength{\parindent}{0pt}
\setlength{\tabcolsep}{0pt}
\raggedbottom

% Section titles: bold black capitals with a thin rule, no colour.
\titleformat{\section}{\large\bfseries}{}{0pt}{\MakeUppercase}[\titlerule]
\titlespacing*{\section}{0pt}{7pt}{5pt}

\setlist[itemize]{leftmargin=1.2em,itemsep=0.5pt,topsep=2pt,parsep=0pt,partopsep=0pt,label=\textbullet}

% \entry{title}{dates}{organisation}{place}
\newcommand{\entry}[4]{%
  \textbf{#1}\hfill #2\par
  \textit{#3}\hfill \textit{#4}\par}
% \project{title}{stack}
\newcommand{\project}[2]{\textbf{#1} \textbar{} \textit{#2}\par}
\newcommand{\skill}[2]{\textbf{#1 :} & #2 \\}

\begin{document}
\small

\begin{center}
  {\fontsize{21}{24}\selectfont\bfseries Amine Mokhtar}\\[4pt]
  {\large Élève ingénieur Cloud \& DevOps}\\[4pt]
  Tunis, Tunisie \textbar{} +216 56 623 537 \textbar{} \href{mailto:aminemokhtar2003@gmail.com}{aminemokhtar2003@gmail.com}\\[2pt]
  \href{https://www.linkedin.com/in/amine-mokhtar-679038277/}{\textbf{LinkedIn}} \textbar{}
  \href{https://github.com/the-amine-mokhtar}{\textbf{GitHub}} \textbar{}
  \href{https://the-amine-mokhtar.github.io/the-amine-mokhtar-the-amine-mokhtar.github.io/}{\textbf{Portfolio}}
\end{center}
\vspace{-6pt}

\section{Profil}
Élève ingénieur en dernière année (Cloud \& Infrastructure) à ESPRIT, à la recherche d'un stage de fin d'études (PFE) en Cloud / DevOps. Expérience pratique en Kubernetes, Docker, Terraform, Ansible, CI/CD et observabilité, dont un stage DevOps chez EY Tunisie.

\section{Expérience professionnelle}
\entry{Stagiaire DevOps}{Juil. 2026 -- Août 2026}{EY Tunisie}{Tunis, Tunisie (hybride)}
\begin{itemize}
  \item Conception d'une infrastructure DevOps auto-hébergée qui détecte les défaillances et récupère sans intervention manuelle.
  \item Provisionnement de 4 VM sur VMware ESXi (Terraform, Ansible) et déploiement d'une application Spring Boot, Angular et MySQL sur un cluster Kubernetes léger.
  \item Mise en place d'un pipeline CI/CD GitHub Actions avec déploiement automatisé et rollback automatique en cas d'échec.
  \item Supervision et auto-remédiation avec Prometheus, Grafana, Alertmanager et Python ; validation par tests de chaos sur 6 scénarios (crash de pods, panne de nœud, déploiement défectueux, perte complète de VM).
\end{itemize}
\vspace{2pt}
\entry{Stagiaire SAP ABAP}{Août 2025 -- Sep. 2025}{SIRYOS}{Tunis, Tunisie}
\begin{itemize}
  \item Développement d'un Business Object de gestion des voyages et réservations avec SAP ABAP RESTful Application Programming (RAP).
  \item Exposition des services via OData et interface utilisateur avec SAP Fiori Elements.
\end{itemize}
\vspace{2pt}
\entry{Stagiaire -- Design graphique}{Juin 2023}{RHIS Software}{Tunis, Tunisie}

\section{Projets}
\project{Infrastructure DevOps Auto-Réparatrice}{VMware ESXi, Kubernetes, Docker, Ansible, Terraform, GitHub Actions, Prometheus, Grafana}
\begin{itemize}
  \item Plateforme entièrement auto-hébergée, provisionnée avec Terraform et configurée avec Ansible, avec remédiation et récupération pilotées par CI/CD.
  \item Tableaux de bord Prometheus et Grafana pour l'observabilité de l'infrastructure et des applications.
\end{itemize}
\vspace{2pt}
\project{CluVerse -- Plateforme SaaS de gestion de clubs étudiants}{OpenStack, Kubernetes, Docker, Ansible, Spring Boot, Angular, Prometheus, Grafana}
\begin{itemize}
  \item Conception et déploiement d'un cloud privé OpenStack distribué avec orchestration Kubernetes des services de la plateforme.
  \item Automatisation du déploiement de l'infrastructure avec Ansible ; supervision de l'application avec Prometheus et Grafana.
\end{itemize}
\vspace{2pt}
\project{MatchUPZ -- Analyse de performances de joueurs}{Python, NLP, SQL, Java, Symfony}
\begin{itemize}
  \item Application web et desktop d'analyse en temps réel des performances de joueurs, avec traitement NLP en Python.
\end{itemize}

\section{Formation}
\entry{Diplôme d'ingénieur en Génie Informatique -- Cloud \& Infrastructure}{Sep. 2022 -- Sep. 2027}{ESPRIT -- École Supérieure Privée d'Ingénierie et de Technologie}{Ariana, Tunisie}

\section{Compétences techniques}
\begin{tabular}{@{}p{5.2cm}l@{}}
  \skill{Cloud \& Virtualisation}{AWS, OpenStack, VMware ESXi}
  \skill{Conteneurs \& Orchestration}{Docker, Kubernetes}
  \skill{IaC \& CI/CD}{Terraform, Ansible, GitHub Actions, Git}
  \skill{Supervision \& Systèmes}{Prometheus, Grafana, Alertmanager, Linux, Bash, Python}
  \skill{Développement}{Java (Spring Boot), JavaScript (Angular, Node.js), PHP (Symfony), C\#, C++}
  \skill{Données \& ERP}{MySQL, Oracle, SQL, PL/SQL, SAP ABAP (RAP, OData, Fiori Elements)}
\end{tabular}

\section{Certifications}
\begin{tabular*}{\textwidth}{@{}l@{\extracolsep{\fill}}r@{}}
  \textbullet\hspace{0.5em}\href{https://www.credly.com/badges/01695568-d031-4853-ba78-202626dd0f25/public_url}{AWS Cloud Operations -- Badge numérique (AWS)} &
  \href{https://www.credly.com/badges/703f6754-2379-41f8-a2a4-a56140597f36}{AWS Security Foundations -- Badge numérique (AWS)} \\
  \textbullet\hspace{0.5em}\href{https://learn.nvidia.com/certificates?j2Df5LzUQ5SHZbV4YqnMKg}{Building RAG Agents with LLMs -- NVIDIA} &
  \href{https://www.efset.org/cert/kRLCmo}{EF SET English Certificate -- 75/100 (C2)} \\
\end{tabular*}

\section{Langues \& activités}
\textbf{Langues :} Arabe (natif), Français (B2), Anglais (C2 -- EF SET)\par
\textbf{Associatif :} IEEE Student Branch ESPRIT, Lions Club Tunis Esprit, Deepflow AI Club

\end{document}
